% =============================================================================
% Presentación — Inteligencia Artificial en Diagnóstico Médico
% Beamer (15 minutos + 5 min preguntas)
% =============================================================================
\documentclass[aspectratio=169,12pt]{beamer}

% --- Theme ---
\usetheme{Madrid}
\usecolortheme{whale}
\setbeamertemplate{navigation symbols}{}
\setbeamertemplate{footline}{%
    \leavevmode%
    \hbox{%
        \begin{beamercolorbox}[wd=.33\paperwidth,ht=2.5ex,dp=1ex,center]{author in head/foot}%
            \usebeamerfont{author in head/foot}\insertshortauthor
        \end{beamercolorbox}%
        \begin{beamercolorbox}[wd=.34\paperwidth,ht=2.5ex,dp=1ex,center]{title in head/foot}%
            \usebeamerfont{title in head/foot}\insertshorttitle
        \end{beamercolorbox}%
        \begin{beamercolorbox}[wd=.33\paperwidth,ht=2.5ex,dp=1ex,right]{date in head/foot}%
            \usebeamerfont{date in head/foot}\insertframenumber{} / \inserttotalframenumber\hspace*{2ex}
        \end{beamercolorbox}%
    }%
}

% --- Encoding & Language ---
\usepackage[utf8]{inputenc}
\usepackage[T1]{fontenc}
\usepackage[spanish]{babel}

% --- Packages ---
\usepackage{graphicx}
\usepackage{booktabs}
\usepackage{tabularx}
\usepackage{amsmath}
\usepackage{xcolor}

% --- Metadata ---
\title[IA en Diagnóstico Médico]{Inteligencia Artificial en Diagnóstico Médico}
\subtitle{Proyecto de Investigación --- Inteligencia Artificial}
\author[A. Segura]{Angel Stward Segura Méndez}
\institute[UNA]{Universidad Nacional\\Sede Regional Brunca --- Campus Pérez Zeledón}
\date{Marzo, 2026}

% =============================================================================
\begin{document}

% ------ PORTADA ------
\begin{frame}
    \titlepage
\end{frame}

% ------ CONTENIDO ------
\begin{frame}{Contenido}
    \tableofcontents
\end{frame}

% =====================================================================
\section{Introducción}
% =====================================================================

\begin{frame}{Introducción}
    % TODO: Contextualización del tema
    \begin{itemize}
        \item Contextualización del tema
        \item Justificación técnica
        \item Relevancia científica e industrial
        \item Objetivo de la investigación
    \end{itemize}
\end{frame}

% =====================================================================
\section{Marco Conceptual}
% =====================================================================

\begin{frame}{Marco Conceptual --- Definiciones}
    % TODO: Definiciones clave
    \begin{itemize}
        \item Inteligencia Artificial
        \item Aprendizaje Automático / Deep Learning
        \item Diagnóstico Asistido por Computadora (CAD)
    \end{itemize}
\end{frame}

\begin{frame}{Marco Conceptual --- Modelos y Arquitecturas}
    % TODO: Modelos relevantes
    \begin{itemize}
        \item Redes Neuronales Convolucionales (CNN)
        \item Arquitecturas relevantes (ResNet, U-Net, etc.)
        \item Transformers en imágenes médicas
    \end{itemize}
\end{frame}

% =====================================================================
\section{Estado del Arte}
% =====================================================================

\begin{frame}{Estado del Arte --- Artículos Analizados}
    % TODO: Resumen de artículos consultados
    \begin{enumerate}
        \item Artículo 1: \textit{Título} --- Resumen breve
        \item Artículo 2: \textit{Título} --- Resumen breve
        \item Artículo 3: \textit{Título} --- Resumen breve
    \end{enumerate}
\end{frame}

\begin{frame}{Estado del Arte --- Comparación de Enfoques}
    % TODO: Tabla comparativa
    \begin{table}
        \centering
        \small
        \begin{tabular}{@{}lccc@{}}
            \toprule
            \textbf{Criterio} & \textbf{Art. 1} & \textbf{Art. 2} & \textbf{Art. 3} \\
            \midrule
            Modelo       & --- & --- & --- \\
            Dataset      & --- & --- & --- \\
            Precisión    & --- & --- & --- \\
            Limitaciones & --- & --- & --- \\
            \bottomrule
        \end{tabular}
    \end{table}
\end{frame}

\begin{frame}{Estado del Arte --- Tendencias Actuales}
    % TODO: Tendencias recientes
    \begin{itemize}
        \item Tendencia 1
        \item Tendencia 2
        \item Tendencia 3
    \end{itemize}
\end{frame}

% =====================================================================
\section{Aplicaciones Reales}
% =====================================================================

\begin{frame}{Aplicaciones Reales}
    % TODO: Casos concretos de implementación
    \begin{itemize}
        \item \textbf{Caso 1:} Descripción breve
        \item \textbf{Caso 2:} Descripción breve
        \item \textbf{Caso 3:} Descripción breve
    \end{itemize}
    \vspace{0.5cm}
    % TODO: Sectores e impacto práctico
\end{frame}

% =====================================================================
\section{Limitaciones y Desafíos Éticos}
% =====================================================================

\begin{frame}{Limitaciones y Desafíos Éticos}
    % TODO: Riesgos, sesgo, escalabilidad, regulación
    \begin{columns}[T]
        \begin{column}{0.48\textwidth}
            \textbf{Limitaciones Técnicas}
            \begin{itemize}
                \item Riesgos técnicos
                \item Sesgo algorítmico
                \item Escalabilidad
            \end{itemize}
        \end{column}
        \begin{column}{0.48\textwidth}
            \textbf{Desafíos Éticos}
            \begin{itemize}
                \item Privacidad de datos
                \item Regulación (FDA/CE)
                \item Responsabilidad legal
            \end{itemize}
        \end{column}
    \end{columns}
\end{frame}

% =====================================================================
\section{Conclusiones}
% =====================================================================

\begin{frame}{Conclusiones}
    % TODO: Síntesis, valoración crítica, perspectivas futuras
    \begin{itemize}
        \item Síntesis de hallazgos principales
        \item Valoración crítica
        \item Perspectivas futuras del área
    \end{itemize}
\end{frame}

% ------ REFERENCIAS ------
\begin{frame}{Referencias}
    % TODO: Lista de artículos consultados (formato APA 7)
    \small
    \begin{enumerate}
        \item Apellido, N. (Año). \textit{Título del artículo 1}. Revista, vol(num), pp. DOI
        \item Apellido, N. (Año). \textit{Título del artículo 2}. Revista, vol(num), pp. DOI
        \item Apellido, N. (Año). \textit{Título del artículo 3}. Revista, vol(num), pp. DOI
    \end{enumerate}
\end{frame}

% ------ CIERRE ------
\begin{frame}
    \centering
    \vspace{1cm}
    {\Huge\textbf{¿Preguntas?}}\\[1cm]
    {\large Angel Stward Segura Méndez}\\[0.3cm]
    {\small Universidad Nacional --- Campus Pérez Zeledón}\\[0.3cm]
    {\small Marzo, 2026}
\end{frame}

\end{document}
